% ================================================================= CHỒNG LẤN VÀ 60 BẢN SẠCH
\subsubsection{Chồng lấn với dữ liệu huấn luyện và 60 bản sạch}
\label{sec:cinc60}

Mục này gom toàn bộ bằng chứng về việc tập đánh giá xuyên hệ ghi của nhóm chứa sẵn dữ liệu huấn luyện,
và trình bày bộ số hiện hành thay cho mọi con số 75 bản ghi đã rút. Mọi con số dưới đây lấy từ
\texttt{benchmark\_dpss/eval\_cinc60\_sach.json} và \texttt{analysis/dulieu\_results.json}; bootstrap
theo bản ghi, 10.000 lần lấy lại, seed 0.

\paragraph{Điều y văn đã ghi nhận --- và điều nhóm bỏ sót.}
Việc tập set-a của PhysioNet/CinC 2013 \textbf{không độc lập} với ADFECGDB \textbf{không phải phát hiện
của nhóm}. Ban tổ chức đã ghi nhận nó ngay từ đầu:

\begin{itemize}[leftmargin=1.6em,itemsep=3pt]
\item Silva và cộng sự, \textit{Computing in Cardiology} 2013;40:149--152, \textbf{Bảng 1} liệt kê
nguồn của 447 bản ghi Challenge, trong đó có dòng ``Abdominal and Direct Fetal ECG Database --- 25''.
\item Clifford và cộng sự, \textit{Physiological Measurement} 2014;35:1521
(DOI 10.1088/0967-3334/35/8/1521), \textbf{Bảng 2}, kèm một câu cảnh báo trực tiếp về một đội dự thi
huấn luyện trên chính bộ đó: bài viết rằng điều này ``\textit{may have led to a bias in the results as
this database was included in set-a, set-b (and possibly a few records in set-c)}''.
\end{itemize}

Câu cảnh báo ấy nằm \textbf{ngay trong bảng ghi chú đọc bài của chính nhóm} --- Bảng
\ref{tab:baihoc}, mục \textbf{p13} (Behar 2014) và mục \textbf{p20} (Clifford 2014), cả hai đều đã tự
viết ra dòng ``việc phải làm ngay là kiểm tra chồng lấn giữa ADFECGDB và CinC 2013''. Nhóm
\textbf{bỏ sót} việc đó suốt các bản 3.2--3.4 và chỉ thực hiện ở vòng 7. Đây là lần đoán sai thứ ba
được phân tích ở \S\ref{sec:doansai}.

\paragraph{Phần thật sự là của nhóm: định danh đúng 15 bản ghi.}
Không nguồn nào nhóm truy cập được cho biết \emph{bản ghi set-a nào} đến từ ADFECGDB. Nhóm xác định
bằng đo lường: tương quan chéo chuẩn hoá (NCC) tính qua FFT trên tín hiệu \textbf{thô chưa lọc}, hạ mẫu
1000 $\to$ 250\,Hz, quét \textbf{mọi} độ trễ, so 75 bản CinC $\times$ 4 kênh với 5 bản ADFECGDB
$\times$ 5 kênh. Hai cài đặt độc lập (\texttt{analysis/dulieu\_audit.py} và
\texttt{analysis/dulieu\_m4b\_verify.py}, viết lại từ đầu, không dùng chung hàm nào) cho cùng một kết
quả: \textbf{đúng 15 bản}, cùng bản gốc, cùng cửa sổ.

Cấu trúc rất đều: mỗi bản ADFECGDB dài 300\,s đóng góp \textbf{đúng ba cửa sổ} 0--60, 120--180 và
240--300\,s; và không phải một kênh mà \textbf{cả bốn kênh} của mỗi bản CinC rò rỉ đều là bản sao
NCC $= 1{,}0000$ của Abdomen\_1..4 \textbf{đúng thứ tự}. Kênh điện cực da đầu không có mặt trong CinC.

\begin{table}[H]\centering\small
\caption{Ánh xạ 15 bản ghi CinC 2013 set-a về 5 bản ghi ADFECGDB dùng huấn luyện. Mọi ô đều có
NCC $= 1{,}0000$ trên cả bốn kênh đúng thứ tự và $|\Delta$RR$|$ trung vị $= 0{,}0$\,ms. Nguồn:
\texttt{analysis/dulieu\_m4b\_verify.json}.}
\label{tab:chonglan15}
\begin{tabular}{@{}L{3.4cm} C{3.0cm} C{3.0cm} C{3.0cm}@{}}
\toprule
\textbf{Bản ADFECGDB (300\,s)} & \textbf{Cửa sổ 0--60\,s} & \textbf{Cửa sổ 120--180\,s} &
\textbf{Cửa sổ 240--300\,s} \\
\midrule
r01 & \xau{a05} & \xau{a04} & \xau{a22} \\
r04 & \xau{a25} & \xau{a20} & \xau{a13} \\
r07 & \xau{a19} & \xau{a23} & \xau{a24} \\
r08 & \xau{a17} & \xau{a15} & \xau{a08} \\
r10 & \xau{a03} & \xau{a14} & \xau{a12} \\
\bottomrule
\end{tabular}
\end{table}

\paragraph{Đối chứng dương và đối chứng âm.}
Một thủ tục dò trùng lặp chỉ đáng tin nếu nó bắt được cái đã biết và không bắt nhầm cái không có.

\begin{itemize}[leftmargin=1.6em,itemsep=3pt]
\item \textbf{Đối chứng dương (cặp đã biết trước).} Cùng thủ tục, chạy trên 5 cặp Silesia B2 ---
ADFECGDB mà nhóm đã xác định ở \S\ref{sec:silesia}, cho NCC 0,856--0,984: B2\_01 $\sim$ r01 (0,860),
B2\_02 $\sim$ r10 (0,984), B2\_07 $\sim$ r04 (0,956), B2\_10 $\sim$ r07 (0,959), B2\_11 $\sim$ r08
(0,856). Thủ tục bắt đúng cái nó phải bắt.
\item \textbf{Đối chứng âm.} 60 bản ghi CinC còn lại có $|$NCC$|$ tối đa chỉ \textbf{0,62} --- cách rất
xa ngưỡng bản sao. Không có vùng xám giữa 0,62 và 1,0000.
\item \textbf{Đối chứng dương thứ hai, bằng hiệu năng.} Trên đúng 15 bản rò rỉ, mô hình 22 ca đạt
\xau{99,87} (trung vị 100; 15/15 bản $\geq 90$; không bản nào dưới 50) --- tức gần bằng hiệu năng
\emph{trong miền} trên chính r01--r10. Trên 60 bản sạch, cùng mô hình, cùng quy tắc, cùng bộ chấm:
\tot{74,28}. Chênh lệch 25 điểm giữa hai tập con của cùng một ``bộ dữ liệu ngoài miền'' là chữ ký của
rò rỉ, không phải của độ khó.
\end{itemize}

\textbf{Mức thổi phồng} khi tính trung bình trên 75 bản thay vì 60 bản sạch, đo riêng cho từng mô hình
và từng quy tắc: mô hình 5 ca $+7{,}18$ điểm; mô hình 12 ca $+6{,}41$; mô hình 22 ca $+5{,}12$; quy tắc
oracle $+3{,}27$. Dải \textbf{3,27--7,18 điểm F1}. Mức thổi phồng lớn nhất rơi vào mô hình yếu nhất ---
đúng như dự đoán, vì 15 bản rò rỉ đều gần 100 với mọi mô hình nên chúng kéo mô hình yếu lên nhiều hơn.
Hệ quả phụ nguy hiểm hơn: chúng \textbf{thu hẹp giả tạo} khoảng cách giữa các mô hình và giữa các quy
tắc chọn kênh, tức làm mọi hiệu số nhóm quan tâm \emph{nhỏ đi} và mọi kết luận ``không khác biệt''
\emph{dễ đạt hơn}.

\paragraph{Bộ số hiện hành: thêm dữ liệu, đo trên 60 bản sạch.}

\begin{table}[H]\centering\small
\caption{PhysioNet/CinC 2013 set-a, \textbf{60 bản ghi sạch}, zero-shot. Hiệu số là mô hình 22 ca trừ
mô hình 5 ca; KTC 95\% bootstrap theo bản ghi 10.000 lần, seed 0; T/Th/H là thắng/thua/hoà theo bản
ghi. Hàng in đậm là \textbf{số chính} của đề tài.}
\label{tab:cinc60}
\begin{tabular}{@{}L{3.6cm} C{1.2cm} C{1.2cm} C{1.2cm} C{1.4cm} C{2.6cm} C{1.5cm}@{}}
\toprule
\textbf{Quy tắc chọn kênh} & \textbf{5 ca} & \textbf{12 ca} & \textbf{22 ca} & \textbf{Hiệu 22--5} &
\textbf{KTC 95\%} & \textbf{T/Th/H} \\
\midrule
\textbf{PSD mù nhãn (chính)} & 64,04 & 67,93 & \tot{74,28} & $+10{,}24$ & $[+7{,}13;\ +13{,}60]$ & 52/3/5 \\
Trung bình 4 kênh & 56,00 & --- & 67,69 & $+11{,}69$ & $[+9{,}03;\ +14{,}37]$ & 58/2/0 \\
Kênh 0 cố định (hậu kiểm, đã rút) & 48,43 & --- & \xau{61,78} & $+13{,}35$ & $[+9{,}67;\ +17{,}10]$ & 50/8/2 \\
Oracle (dùng nhãn, chỉ là trần) & 72,76 & --- & \textit{83,60} & $+10{,}84$ & $[+7{,}55;\ +14{,}37]$ & 49/3/8 \\
\midrule
\multicolumn{7}{@{}l}{\textit{Biến thể 53 bản (loại thêm 7 bản chú thích sai), quy tắc PSD:}} \\
PSD mù nhãn, $n = 53$ & 64,19 & --- & \tot{75,27} & $+11{,}08$ & $[+7{,}54;\ +14{,}82]$ & --- \\
\bottomrule
\end{tabular}
\end{table}

Trị số $p$ Wilcoxon ghép cặp cho hiệu số 22 ca trừ 5 ca: $3{,}5 \cdot 10^{-10}$ (PSD),
$2{,}7 \cdot 10^{-11}$ (trung bình 4 kênh), $4{,}9 \cdot 10^{-9}$ (kênh 0), $1{,}1 \cdot 10^{-9}$
(oracle). Với mô hình 22 ca và quy tắc PSD: \textbf{33/60} bản đạt $\geq 90$, \textbf{16/60} bản còn
dưới 50, trung vị 95,07. Đường 5 $\to$ 12 $\to$ 22 ca (64,04 $\to$ 67,93 $\to$ 74,28) cho thấy lợi ích
của dữ liệu \textbf{chưa bão hoà}: riêng bước 12 $\to$ 22 còn đáng $+6{,}35$ điểm
$[+4{,}04;\ +8{,}95]$. Power-MF đơn kênh trên đúng 60 bản này đạt \textbf{55,97}; chênh lệch so với
\hethong{} là $+18{,}32$ $[+13{,}35;\ +23{,}64]$, thắng 49 thua 5 hoà 6 (\S\ref{sec:powermf}).

\paragraph{Quy tắc chọn kênh: bảy phương án, một quy tắc chỉ định trước, và một kết quả trượt.}
Bảy phương án mù nhãn thay cho quy tắc PSD được ghi vào một bản khai báo
(\texttt{analysis/chonkenh\_khaibao\_truoc.json}) \textbf{trước khi} bất kỳ phương án nào được chấm
trên CinC, kèm quy tắc quyết định ``quy tắc chỉ định trước $=$ phương án có F1 trong miền cao nhất trên
22 chủ thể'' --- quy tắc đó chỉ ra \texttt{gate}. Bản khai báo được viết \emph{sau} khi đã có F1 từng
kênh và \textbf{không được neo bởi bên thứ ba}, nên nhóm gọi nó là \textbf{bản khai báo trước} và
\textbf{không} gọi là ``tiền đăng ký''.

\begin{table}[H]\centering\footnotesize
\caption{Bảy quy tắc chọn kênh mù nhãn trên \textbf{60 bản sạch}, mô hình 22 ca. Hiệu số so với quy tắc
PSD, KTC 95\% bootstrap theo bản ghi, Holm hiệu chỉnh trên họ 7 quy tắc. Cột ``\% dư địa'' là phần
khoảng cách PSD $\to$ oracle (9,32 điểm) mà quy tắc lấy lại được.}
\label{tab:quytac60}
\begin{tabular}{@{}L{2.0cm} L{2.3cm} C{0.9cm} C{1.1cm} C{2.3cm} C{1.2cm} C{1.4cm} C{1.0cm}@{}}
\toprule
\textbf{Quy tắc} & \textbf{Vai trò} & \textbf{F1} & \textbf{Hiệu} & \textbf{KTC 95\%} &
\textbf{$p_{\text{Holm}}$} & \textbf{T/Th/H} & \textbf{\% dư địa} \\
\midrule
\texttt{psd} & tham chiếu & 74,28 & --- & --- & --- & 0/0/60 & 0 \\
\texttt{gate} & \textbf{chỉ định trước} & 80,72 & $+6{,}44$ & $[+2{,}49;\ +11{,}10]$ & \xau{0,051} & 16/7/37 & 69,1 \\
\texttt{gate4} & khai báo trước & 81,01 & $+6{,}72$ & $[+2{,}85;\ +11{,}18]$ & \tot{0,015} & 17/5/38 & 72,2 \\
\texttt{peakprob} & \textbf{hậu kiểm} & 82,01 & $+7{,}73$ & $[+3{,}82;\ +12{,}41]$ & \tot{0,0039} & 19/6/35 & 82,9 \\
\texttt{rrcv} & khai báo trước & 80,00 & $+5{,}72$ & $[+1{,}61;\ +10{,}41]$ & 0,41 & 15/10/35 & 61,3 \\
\texttt{rrplaus} & khai báo trước & 78,87 & $+4{,}59$ & $[+0{,}88;\ +8{,}93]$ & 0,41 & 17/9/34 & 49,3 \\
\texttt{learned} & khai báo trước & 77,80 & $+3{,}52$ & $[-0{,}77;\ +8{,}48]$ & 1,00 & 15/13/32 & 37,8 \\
\texttt{fuse} ($k = 2$) & khai báo trước & 76,46 & $+2{,}18$ & $[-0{,}40;\ +4{,}98]$ & 1,00 & 20/20/20 & 23,4 \\
\midrule
\textit{oracle} & trần (dùng nhãn) & \textit{83,60} & $+9{,}32$ & $[+5{,}10;\ +14{,}15]$ & --- & 25/0/35 & 100 \\
\bottomrule
\end{tabular}
\end{table}

Ba điều phải đọc cùng nhau, và chúng không cùng chiều:

\begin{enumerate}[leftmargin=1.8em,itemsep=4pt]
\item \textbf{Quy tắc chỉ định trước trượt.} \texttt{gate} đạt $+6{,}44$ điểm, KTC 95\% loại trừ 0,
Wilcoxon $p = 0{,}010$ --- nhưng sau hiệu chỉnh Holm trên họ 7 quy tắc,
$p_{\text{Holm}} = \xau{0{,}051}$. Theo đúng quy tắc quyết định đã ghi \emph{trước khi} chạy, đây là
\textbf{trượt}, và nhóm báo cáo nó là trượt. Nhóm \textbf{không} đổi họ kiểm định, \textbf{không} bỏ
bớt quy tắc khỏi họ để hạ ngưỡng, và \textbf{không} viết ``gần đạt ý nghĩa''.
\item \textbf{Quy tắc mạnh nhất lại là quy tắc hậu kiểm.} \texttt{peakprob} --- chọn kênh có xác suất
đỉnh trung bình cao nhất theo chính đầu ra của bộ dò; không tham số, không huấn luyện --- đạt 82,01
($+7{,}73$; $p_{\text{Holm}} = 0{,}0039$), lấy lại 82,9\,\% dư địa tới oracle và hạ số bản dưới 50 từ
16 xuống 9. Nhưng nó \textbf{được chọn sau khi nhìn kết quả}. Theo đúng nguyên tắc nhóm tự đặt ở
\S\ref{sec:doansai}, nó \textbf{chỉ là giả thuyết}, không phải kết quả khẳng định.
\item \textbf{Trong miền, không quy tắc nào tách được khỏi PSD.} Trên 22 chủ thể, F1 bão hoà gần 100
nên nhóm lặp lại so sánh trên \textbf{thang logit} (\texttt{analysis/xacnhan\_results.json}). Ở đó thứ
tự \textbf{đổi}: quy tắc đứng đầu là \texttt{gate} ($+0{,}225$ logit; KTC $[+0{,}004;\ +0{,}589]$;
$p_{\text{Holm}} = 0{,}475$), còn \texttt{peakprob} chỉ được $+0{,}134$ $[-0{,}038;\ +0{,}407]$,
$p_{\text{Holm}} = 1{,}00$. Không quy tắc nào sống sót qua Holm. Nói cách khác: \textbf{bằng chứng
trong miền nghiêng về quy tắc chỉ định trước, bằng chứng ngoài miền nghiêng về quy tắc hậu kiểm} ---
và hai bằng chứng đó không thể cùng được dùng để chốt một quy tắc.
\end{enumerate}

\paragraph{Vì sao không thể đóng câu hỏi này ở vòng 7.}
Cách duy nhất để biến \texttt{peakprob} từ giả thuyết thành kết quả là chấm nó trên một bộ dữ liệu thứ
ba \textbf{có nhãn fQRS thật}. Nhóm đã kiểm bốn khả năng và \textbf{không bộ nào dùng được}:

\begin{table}[H]\centering\small
\caption{Bốn bộ dữ liệu được kiểm làm tập xác nhận thứ ba (vòng 7). Nguồn:
\texttt{analysis/xacnhan\_results.json}, \texttt{analysis/XACNHAN.md}.}
\label{tab:bothuba}
\begin{tabular}{@{}L{2.7cm} L{3.3cm} C{1.8cm} L{6.2cm}@{}}
\toprule
\textbf{Bộ dữ liệu} & \textbf{Có trên đĩa} & \textbf{Nhãn fQRS} & \textbf{Bằng chứng} \\
\midrule
NIFEADB & 26 bản \texttt{.dat}/\texttt{.hea} & \xau{Không} & PhysioNet 1.0.0 không có
\texttt{.qrs}/\texttt{.atr}/\texttt{ANNOTATORS}; cờ \texttt{HAS\_FETAL\_ANNOTATIONS} trong
\texttt{nifeadb\_loader} đặt False \\
NInFEA & 2 \texttt{.dat} $+$ 3 \texttt{.hea} (tải dò) & \xau{Không} & Không có tệp chú thích nhịp;
tham chiếu gốc của bộ này là Doppler \\
nifecgdb (ecgca102) & 1 bản EDF $+$ \texttt{.qrs} & \xau{Không} & 375 nhãn trên 270\,s, RR trung vị
0,695\,s $=$ 86 nhịp/phút $\Rightarrow$ đó là QRS \textbf{mẹ}; ngoài ra 55 bản chỉ từ một sản phụ \\
CinC 2013 set-b & Không & \xau{Không} & Nhãn chưa bao giờ được công bố \\
\bottomrule
\end{tabular}
\end{table}

\begin{canhbao}[title={Phát biểu nhóm dùng cho mục này --- và phát biểu nhóm KHÔNG dùng}]
\textbf{Nhóm viết:} ``Như ban tổ chức đã ghi nhận [Silva 2013, Bảng 1; Clifford 2014, Bảng 2], set-a
chứa bản ghi của ADFECGDB; chúng tôi xác định bằng đo lường \emph{đúng 15 bản nào} và \emph{mức thổi
phồng} 3,27--7,18 điểm F1, rồi chấm lại toàn bộ trên 60 bản còn lại.''

\textbf{Nhóm KHÔNG viết:} ``nhóm phát hiện rò rỉ dữ liệu trong CinC 2013''. Đó là một phát biểu sai về
quyền ưu tiên và đã bị rút (\S\ref{sec:doansai}).

\textbf{Nhóm cũng KHÔNG viết:} ``\texttt{peakprob} nâng F1 xuyên hệ ghi lên 82,01'' như một kết quả.
Con số 82,01 có thật và chạy lại được, nhưng quy tắc sinh ra nó là hậu kiểm, quy tắc chỉ định trước đã
trượt Holm, thang logit trong miền lại xếp hạng khác, và \textbf{không có bộ dữ liệu công khai thứ ba
có nhãn fQRS thật} để phân xử. Trạng thái đúng của \texttt{peakprob} là \textbf{giả thuyết mạnh, chưa
được xác nhận}. Việc xác nhận nó cần \emph{dữ liệu mới}, không cần thêm phân tích trên dữ liệu cũ.
\end{canhbao}
